\underline{Exercise 7.1}
\begin{enumerate}
% Question 3
\item[3.] Find $f^{\prime}(1)$ and $f^{\prime}(2)$ for the following functions:
\begin{multicols}{3}
\begin{enumerate}
\item[(a)] $y=f(x)=18 x$
\item[(b)] $y=f(x)=c x^{3}$
\item[(c)] $f(x)=-5 x^{-2}$
\item[(d)] $f(x)=\frac{3}{4} x^{4 / 3}$
\item[(e)] $f(w)=6 w^{1 / 3}$
\item[(f)] $f(w)=-3 w^{-1 / 6}$
\end{enumerate}
\end{multicols}
\end{enumerate}

\underline{Exercise 7.2}
\begin{enumerate}
% Question 3
\item[3.] Differentiate the following by using the product rule:
\begin{multicols}{2}
\begin{enumerate}
\item[(d)] $(a x-b)\left(c x^{2}\right)$
\item[(e)] $(2-3 x)(1+x)(x+2)$
\end{enumerate}
\end{multicols}

% Question 7
\item[7.] Find the derivatives of:
\begin{multicols}{2}
\begin{enumerate}
\item[(a)] $\left(x^{2}+3\right) / x$
\item[(b)] $(x+9) / x$
\item[(c)] $6 x /(x+5)$
\item[(d)] $\left(a x^{2}+b\right) /(c x+d)$
\end{enumerate}
\end{multicols}

% Question 8
\item[8.] Given the function $f(x)=a x+b$, find the derivatives of:
\begin{multicols}{4}
\begin{enumerate}
\item[(a)] $f(x)$
\item[(b)] $x f(x)$
\item[(c)] $1 / f(x)$
\item[(d)] $f(x) / x$
\end{enumerate}
\end{multicols}
\end{enumerate}

\underline{Exercise 7.3}
\begin{enumerate}
% Question 1
\item[1.] Given $y=u^{3}+2 u$, where $u=5-x^{2}$, find $d y / d x$ by the chain rule.

% Question 2
\item[2.] Given $w=a y^{2}$ and $y=b x^{2}+c x$, find $d w / d x$ by the chain rule.

% Question 3
\item[3.] Use the chain rule to find $d y / d x$ for the following:
\begin{multicols}{3}
\begin{enumerate}
\item[(a)] $y=\left(3 x^{2}-13\right)^{3}$
\item[(b)] $y=\left(7 x^{3}-5\right)^{9}$
\item[(c)] $y=(a x+b)^{5}$
\end{enumerate}
\end{multicols}

% Question 4
\item[4.] Given $y=(16 x+3)^{-2}$, use the chain rule to find $d y / d x$. Then rewrite the function as $y=1 /(16 x+3)^{2}$ and find $d y / d x$ by the quotient rule. Are the answers identical?

% Question 5
\item[5.] Given $y=7 x+21$, find its inverse function. Then find $d y / d x$ and $d x / d y$, and verify the inverse-function rule. Also verify that the graphs of the two functions bear a mirror-image relationship to each other.

% Question 6
\item[6.] Are the following functions strictly monotonic? For each strictly monotonic function, find $d x / d y$ by the inverse-function rule.
\begin{multicols}{2}
\begin{enumerate}
\item[(a)] $y=-x^{6}+5 \quad(x>0)$
\item[(b)] $y=4 x^{5}+x^{3}+3 x$
\end{enumerate}
\end{multicols}
\end{enumerate}
